%% P21 \dash{} Optymalizacja z ograniczeniami i mnożniki Lagrange'a
%%
%% Zalążek. Poniższy opis jest kontraktem tego programu: co ma uzasadnić i co
%% ma dać czytelnikowi. Napisanie programu oznacza usunięcie bloku
%% \programstub{} i zastąpienie go programem. Jeśli po przerobieniu materiału
%% opis okaże się błędny, popraw go i odnotuj sprzeczność w CLAUDE.md w sekcji
%% *Resolved questions*.

\program{Optymalizacja z ograniczeniami i mnożniki Lagrange'a}
\label{prog:P21}

\programstub{%
\textit{[Opis do przetłumaczenia \dash{} patrz wydanie angielskie.]}

Argument: at a constrained optimum the gradients of objective and constraint
are parallel; the multiplier is the exchange rate between them. Payoff: a
Lagrange multiplier is a price \dash{} how much objective you buy per unit of
constraint relaxed \dash{} and that reading turns the \code{beta} in a KL-
penalised objective from a tuning knob into a quantity with units; the
equivalence between a hard KL constraint and a KL penalty, which is why PPO
and DPO are talking about the same problem; projection onto a set as the
other way to enforce a constraint; KKT conditions stated for recognition,
with the inequality case sketched rather than developed.

\medskip
\textbf{Planowana długość:} 50 ramek.
}
